% Part: counterfactuals
% Chapter: introduction
% Section: strict-conditional

\documentclass[../../../include/open-logic-section]{subfiles}

\begin{document}

\olfileid{cnt}{int}{str}

\olsection{Kondisional Ketat}

Lewis ngenalake \emph{kondisional ketat}~$\strictif$ lan ngandharake
yen kondisional iki, dudu kondisional material, cocog karo implikasi.
Ing logika modal aletik, $!A \strictif !B$ bisa ditetepake minangka
$\Box(!A \lif !B)$. Dadi kondisional ketat bener (ing sawijining
donya) yen lan mung yen kondisional material kang cocog iku niscaya.

Kepriye kondisional ketat tumrap paradoks kondisional material?
Kondisional ketat kanthi anteseden kang luput utawa kanthi konsekuen
kang bener bisa bener lan uga bisa luput. Kajaba iku,
$(!A \strictif !B) \lor (!B \strictif !A)$ ora sah. Kondisional
ketat $!A \strictif !B$ uga ora ekuivalen karo $\lnot !A \lor !B$,
mula ora fungsional tumrap bebener.

Ing sistem S5, kita nduweni:
\begin{align}
  !A \strictif !B & \Entails \lnot !A \lor !B
  \text{ nanging:}\\
  \lnot !A \lor !B & \Entails/ !A \strictif !B\\
  !B & \Entails/ !A \strictif !B\\
  \lnot !A & \Entails/ !A \strictif !B\\
  \lnot(!A \strictif !B) & \Entails/ !A \land \lnot !B
  \text{ nanging:}\\
  !A \land \lnot !B & \Entails \lnot(!A \strictif !B)
  \intertext{Nanging kondisional ketat isih ndhukung modus ponens:}
  !A, !A \strictif !B & \Entails !B
  \intertext{Konsekuen-konsekuen saka kondisional ketat kanthi
    anteseden kang padha bisa digandhengake:}
  !A \strictif !B,  !A \strictif !C & \Entails !A \strictif (!B \land !C)
  \intertext{Penguwatan anteseden lumaku kanggo kondisional ketat:}
  !A \strictif !B & \Entails (!A \land !C) \strictif !B
  \intertext{Kondisional ketat uga transitif:}
  !A \strictif !B, !B \strictif !C & \Entails !A \strictif !C
  \intertext{Pungkasan, kondisional ketat ekuivalen karo
    kontrapositife:}
  !A \strictif !B & \Entails \lnot !B \strictif \lnot !A\\
  \lnot !B \strictif \lnot !A & \Entails !A \strictif !B
\end{align}

\begin{prob}
  Wenehana tuladha panyanggah ing \Log{S5} tumrap relasi konsekuensi
  kang ora lumaku kanggo kondisional ketat, yaiku:
  \begin{enumerate}
  \item $\lnot p \Entails/ \Box(p \lif q)$
  \item $q \Entails/ \Box(p \lif q)$
  \item $\lnot\Box(p \lif q) \Entails/ p \land \lnot q$
  \item $\Entails/ \Box(p \lif q) \lor \Box(q \lif p)$
  \end{enumerate}
\end{prob}

\begin{prob}
  Tuduhna yen relasi konsekuensi kang sah lumaku kanggo kondisional
  ketat kanthi menehi bukti ing \Log{S5} kanggo:
  \begin{enumerate}
  \item  $\Box(!A \lif !B) \Entails \lnot !A \lor !B$
  \item $!A \land \lnot !B \Entails \lnot\Box(!A \lif !B)$
  \item $!A, \Box(!A \lif !B) \Entails !B$
  \item $\Box(!A \lif !B), \Box(!A \lif !C) \Entails \Box(!A \lif (!B
    \land !C))$
  \item $\Box(!A \lif !B) \Entails \Box((!A \land !C) \lif !B)$
  \item $\Box(!A \lif !B), \Box(!B \lif !C) \Entails \Box(!A
    \lif !C)$
  \item $\Box(!A \lif !B) \Entails \Box(\lnot !B \lif \lnot !A)$
  \item $\Box(\lnot !B \lif \lnot !A) \Entails \Box(!A \lif !B)$
  \end{enumerate}
  \end{prob}

Nanging kondisional ketat isih nduweni ``paradoks'' dhewe. Kaya
kondisional material kanthi anteseden kang luput utawa konsekuen
kang bener iku bener, kondisional ketat kanthi anteseden kang
\emph{niscaya} luput utawa konsekuen kang niscaya bener iku bener.
Kajaba iku, ing S5 saben kondisional ketat kang bener iku niscaya
bener, lan saben kondisional ketat kang luput iku niscaya luput.
Kanthi tembung liya, kita nduweni
\begin{align}
  \Box \lnot !A & \Entails !A \strictif !B\\
  \Box !B & \Entails !A \strictif !B\\
  !A \strictif !B& \Entails \Box(!A \strictif !B)\\
  \lnot(!A \strictif !B) & \Entails \Box\lnot(!A \strictif !B)
\end{align}
Iki ora nuwuhake masalah yen kita nganggep $\strictif$ minangka
``implies'' (``ngimplikasikake''). Relasi konsekuensi logis iku
fakta matematika, mula ora bisa kontingen. Nanging iki nuwuhake
masalah yen kita arep nggunakake $\strictif$ minangka konektif
logis kang makili ``if \dots then \dots'', mligine rong relasi
pungkasan. Mesthi ana pratelan ``if \dots then \dots'' kang bener
kanthi kontingen utawa luput kanthi kontingen---nyatane, lumrahe
pratelan iku ora niscaya lan uga ora mokal.

\begin{prob}
  Wenehana bukti ing \Log{S5} kanggo:
  \begin{enumerate}
    \item  $\Box \lnot !A \Entails !A \strictif !B$
    \item $!A \strictif !B \Entails \Box(!A \strictif !B)$
    \item $\lnot(!A \strictif !B)  \Entails \Box\lnot(!A \strictif !B)$
  \end{enumerate}
  Gunakna definisi $\strictif$ kanggo iku.
\end{prob}

\end{document}
